%!TEX root = paper.tex

\section{Introduction}
\label{sec:intro}
Randomized sampling is an important tool in the design of efficient algorithms.
A random subset often preserves key properties of the entire data set,
allowing one to run algorithms on a small sample.
A particularly useful instance of this phenomenon is row sampling of matrices.
For a $n \times d$ matrix $\AA$ where $n \gg d$ and any error parameter
$\epsilon > 0$, we can find $\AA'$ with a few (rescaled) rows of $\AA$ such that
\begin{equation*}
\norm{\AA \xx}_{p} \approx_{1 + \epsilon} \norm{\AA' \xx}_{p}
\end{equation*}
for all vectors $\xx \in \Re^{d}$. Here $\approx_{1 + \epsilon}$
denotes a multiplicative error between $(1 + \epsilon)^{-1}$ and $(1 + \epsilon)$.

Originally studied in statistics~\cite{Talagrand90,RudelsonV07,Tropp12},
the row sampling problem has received much attention recently in randomized numerical linear
algebra~\cite{DrineasMM06,DrineasMMW11,ClarksonW13,MahoneyM13,NelsonN12,LiMP13,CohenLMMPS14:arxiv}
and in graph algorithms as spectral sparsification~\cite{SpielmanS08:journal,BatsonSS12,KoutisLP12}.
These works led to a good understanding of row sampling for the $p = 2$ case.
If the rows of $\AA$ are sampled with probabilities proportional to their
statistical leverage scores, matrix Chernoff bounds~\cite{AhlswedeW02,RudelsonV07,Tropp12}
state that $\AA'$ with $O(d \log{d} / \epsilon^{-2})$ is a good approximation with high probability. 
Recently, Clarkson and Woodruff developed oblivious subspace embeddings
that brought the runtime of these algorithms down to input-sparsity time~\cite{ClarksonW13}.
A derandomized procedure by Batson et al.~\cite{BatsonSS12} can also
reduce the number of rows in $\AA'$ to $O(d / \epsilon^2)$, at the cost of a worse,
but still polynomial, runtime.

Substantial progress has also been made for other values of $p$~\cite{DasguptaDHKM09,SohlerW11,ClarksonDMMMW13,MahoneyM13},
leading to input-sparsity time algorithms that return samples
with about $d^{2.5}$ rows when $1 \leq p \leq 2$.
The $p = 1$ case is of particular interest due to its relation
to robust regression, or $\ell_1$-regression~\cite{Candes06}.
In this setting, the existence of samples of size $O(d \log{d})$ was
shown by Talagrand~\cite{Talagrand90} using Banach space theory.

Talagrand's result, and the best known sampling results for general
$\ell_p$ norms~\cite{BourgainLM89,TalagrandLp}, are based on a
``change of density'' construction originally due to Lewis~\cite{lewis}. This construction
assigns a weight, analogous to a leverage score, to each row; these can be used
directly as sampling probabilities.  We will refer to these weights as ``Lewis weights''.
Given their direct use as sampling probabilities, the primary algorithmic challenge
is to be able to compute, or at least approximate, these weights.  This paper provides
the first input-sparsity time algorithms, and in fact the first polynomial time algorithms,
for this problem.  That in turn leads to the first polynomial time algorithmic versions
of the Talagrand and other Lewis weight-based results.

In particular, we give a simple iterative algorithm that approximates Lewis weights
through repeated computations of statistical leverage scores.
Sampling by these approximate weights then leads to the following result:

\begin{theorem}
\label{thm:main}
Given a matrix $\AA$ and any error parameter $\epsilon > 0$,
there is a distribution of matrices $\SS$ with $O(d \log d \epsilon^{-2})$
rows and one nonzero entry per row such that with high probability
\begin{equation*}
\norm{\SS \AA \xx}_1 \approx_{1 + \epsilon} \norm{\AA \xx}_1
\qquad \text{for all } \xx \in \Re_{d}.
\end{equation*}
Furthermore, we can sample from this distribution using $O(\log \log{n})$ calls
to computing 2-approximate statistical leverage scores of matrices of the
form $\WW \AA$ where $\WW$ is a non-negative diagonal matrix.
\end{theorem}
This routine can be combined with any algorithm for computing (approximate) statistical leverage scores.
By invoking input-sparsity time routines~\cite{ClarksonW13,MahoneyM13,
NelsonN12,LiMP13,CohenLMMPS14:arxiv} and approximating intermediate
objects to a coarser granularity, we obtain algorithms that compute $\AA'$ with
$O(d \log{d} \epsilon^{-2})$ rows in $O(\nnz(\AA) + d^{\omega + \theta} )$ time.
Here $\omega$ is the matrix-multiplication exponent and $\theta > 0$ is any constant.
This is a substantial improvement over previous results.
The previous best algorithm with a comparable running time gives a sample size of
about $d^{3.66}$~\cite{LiMP13}; algorithms for obtaining samples with
$O(d^{2.5})$ rows rely on an expensive ellipsoidal rounding algorithm~\cite{DasguptaDHKM09},
and therefore have prohibitively large $\poly(d)$ terms.

We also give input-sparsity time algorithms for sampling for general $\ell_p$ norms.
However, the sample count is slightly higher for $p \in (1,2)$ and substantially higher
(around $d^{p/2}$) for $p > 2$.  Furthermore, the runtime gets substantially worse as
$p$ approaches and exceeds 4 (we  use the ellipsoid algorithm in this regime).
Nonetheless, these algorithms do much better than previous ones: they give a near-optimal
sample count, and their dependence on $d$ (though high) is of the form
$d^{p/2 (1 + \theta) + C}$
(better than could be obtained in approaches based on resparsifying the
outputs of existing algorithms).

A table summarizing our main algorithmic results is in Figure~\ref{fig:summary}.
The bounds omit logarithmic factors in the runtime, $d^{\theta}$ tradeoffs
with input sparsity terms, success probabilities, and $\epsilon$
and constant factor $p$ dependences.
All of these algorithms can be run with an $\nnz(\AA) \log n$ term and also
support a tradeoff to get pure input sparsity time.


\begin{figure}

\begin{center}

\begin{tabular}{| l | l | l | l | }
\hline
$p$ & polynomial runtime term & Sample count & Relevant Sections\\ \hline
$p = 1$ & $d^{\omega}$ & $d \log d$ &
Sections~\ref{sec:overview},~\ref{sec:compute},~\ref{sec:ellp},~and~\ref{sec:concentration}\\ \hline
$1 < p < 2$ & $d^{\omega}$ & $d \log d \log \log^2 d$ & Sections~\ref{sec:overview},~\ref{sec:compute}, and~\ref{sec:ellp}\\ \hline
$2 < p < 4$ & $d^{\omega}$ & $d^{p/2} \log d$ &
Sections~\ref{sec:overview},~\ref{sec:compute},~and~\ref{sec:ellp}\\ \hline
$2 \leq p $ & $d^{p/2+C}$ & $d^{p/2} \log d$ &
Sections~\ref{sec:overview},~\ref{sec:convex},~\ref{sec:sparseconvex}~and~\ref{sec:ellp}\\ \hline
\end{tabular}

\end{center}

\caption{Summary of Our Algorithmic Results}
\label{fig:summary}
\end{figure}


In Section~\ref{sec:ellp}, we briefly go over known statistical results~\cite{Talagrand90,
TalagrandLp,BourgainLM89} on the concentration of sampling by Lewis weights.
One issue with invoking these results is that the bounds in them do not explicitly
describe non-uniform sampling based on upper bounds of Lewis weights.
Instead, the more common form of a main theorem in these results is that
if the Lewis weights are uniformly small, uniformly sampling half the rows gives
a good approximation.
We use standard techniques (which we describe in more detail in
Appendix~\ref{sec:reduct}) to show that the latter implies the former.

Additionally, in the case of $\ell_1$, we give an alternate, elementary proof
of the concentration of the sampling process in Section~\ref{sec:concentration}.
It simplifies many components of previous
proofs~\cite{Talagrand90,Pisier99:book} while following the same basic strategy.
The bound is slightly weaker in some artificial parameter ranges, but is also stronger
in giving much sharper tails.  It is arguably even simpler than proofs of $\ell_2$ matrix Chernoff
bounds~\cite{Vershynin09,Harvey11}.  Unfortunately, the picture for general $\ell_p$ seems
to be much more difficult: published proofs for $p \not \in \{ 1, 2 \}$ use deep results from
Gaussian processes.  We consider it an open problem to find elementary proofs for these ranges.

Our results significantly improve algorithms for the well studied
$\ell_p$, and specifically $\ell_1$, row sampling problem.
We also give a version of $\ell_1$ matrix concentration bound that's
analogous to the widely-used $\ell_2$ matrix concentration bounds.
We believe these results and the simplicity of our techniques
show the usefulness of Lewis weights as an algorithmic tool
for randomized numerical linear algebra in $p$-norms.

The paper is structured as follows: in Section~\ref{sec:overview} we formalize
the matrix row sampling problem and give an overview of our result.
Our simple iterative algorithm for computing approximate Lewis weights is in Section~\ref{sec:compute};
this applies to all $p < 4$.
An alternative approach to computing Lewis weights based on convex optimization, valid for all $p \geq 2$,
is in Section~\ref{sec:convex}.
Section~\ref{sec:properties} gives properties of Lewis weights useful in some of our arguments.
An input-sparsity time algorithm using the convex optimization approach is in Section~\ref{sec:sparseconvex}.
Section~\ref{sec:ellp} summarizes the known sampling results for each range of $p$, and describes
a way to take existing proofs and obtain analyses of our simple sampling procedure.
Section~\ref{sec:concentration} gives most of our elementary proof of the validity
of $\ell_1$ sampling by Lewis weight.
Finally, Appendix~\ref{sec:properties-proofs} proves the properties from Section~\ref{sec:properties},
and Appendix~\ref{sec:reduct} gives the proof of the reduction from our sampling procedure
given in Section~\ref{sec:ellp}.
